\paragraph{Why pullback iteration does not improve an existing primitive factor.}
Let $X$ be a smooth projective geometrically connected curve
of genus at least two. Let $D\to X$ be a geometrically connected
etale cover of odd degree $n>1$ prime to $p$, geometrically
Galois with cyclic group, and let $a:Y\to X$ be a finite etale
cover of degree $d$ prime to $p$.
Assume that all curves and maps are defined over the same finite
field $k$ of characteristic $p$, and that
\[
                         D'=D\times_XY
 \tag{PI.1}
\]
is geometrically connected. Then $D'\to Y$ has the same
geometric cyclic degree $n$. Let $P_n$ and $P'_n$ be the
actual primitive conductor-$n$ subvarieties in $J_D,J_{D'}$.
For every finite extension $k'/k$, pullback gives an injection
\[
             P_n(k')[p^\infty]\lhook\joinrel\longrightarrow
                              P'_n(k')[p^\infty],
 \qquad
             f(P_n)\le f(P'_n),
 \tag{PI.2}
\]
where $f$ denotes geometric $p$-rank. In particular the old
rational exponent divides the new one.

\begin{proof}
Let $h:D'\to D$ be the projection, of degree $d$.
The deck groups are naturally identified by base change;
write their corresponding generators as $\sigma,\sigma'$.
Pullback and norm are equivariant for these generators.
For the canonical integral projectors $p_n=ne_n$ and
$p'_n=ne'_n$ of (IS.2)--(IS.3), this gives
\[
 h^*p_n=p'_nh^*,\qquad h_*p'_n=p_nh_*.
 \tag{PI.3}
\]
Consequently $h^*$ maps the actual image $P_n$ into the actual
image $P'_n$, and $h_*$ maps $P'_n$ back into $P_n$.
These maps are defined over $k$, even if the geometric deck
groups themselves are not constant. Their composite on $P_n$ is
\[
                          h_*h^*=[d].
 \tag{PI.4}
\]
Since $p\nmid d$, multiplication by $d$ is invertible on every
$p$-primary torsion group. Thus $h^*$ is injective on those
groups over $\bar k$ and over every $k'$. Restricting to
geometric $p$-torsion proves the $p$-rank inequality, and
the rational injections prove the exponent divisibility.
No arbitrary isogenous representative of a primitive factor
is substituted.
\end{proof}

Thus a nonzero old rational primitive $p$-primary group cannot
become zero by this construction. If the same-$n$ numerical
inequality of A47 already fails for the old exponent, it still
fails for the new exponent. An additional constant extension
cannot improve either conclusion, since old rational torsion
remains rational. This concerns the vanishing and exponent
inputs to A46/A47; it does not forbid a newly chosen cyclic
cover of $Y$ from satisfying them.

Connectedness and retention of the conductor are essential.
If the fibre product is disconnected and one selects a
component, the cyclic degree or conductor can drop; no
same-conductor assertion is made in that situation.
The pointing also remains a separate hypothesis when applying
a section criterion: rational lifts in both given covers
produce a rational point of the fibre product, but the
injection alone does not supply those lifts.

\paragraph{Raynaud's simultaneous theta incidences are not nilpotence.}
Raynaud's Theorem~2, printed p.~76, gives a prime-to-$p$
solvable Galois cover $Y\to X$ and a finite representation
factoring through its group which has no theta divisor.
This is a geometric existence statement over an algebraic
closure; finite-field descent and pointing may require a
finite extension of the initial field.
Proposition~1(5), printed p.~75, then gives
\[
 a^{(1)*}J_{X^{(1)}}\ \subset\ \operatorname{Supp}(\Theta_Y),
 \tag{PI.5}
\]
where the theta divisor is that of locally exact differentials
\cite[Proposition~1(5), Theorem~2]{RaynaudNonordinary2000}.
In particular all the multiples of a pulled-back torsion class
lie on this theta divisor, provided the class retains its order.
The theorem and its construction are stronger than a single
nonordinary-cover assertion.

Nevertheless, if an original cyclic cover has positive
primitive $p$-rank, (PI.2) makes the connected pulled-back
cover's primitive $p$-rank positive as well. The simultaneous
incidences therefore cannot be interpreted as a statement that
the primitive Frobenius operator is nilpotent. In genus
$g(Y)>2$, a nontrivial character has coherent cohomology
dimension $g(Y)-1>1$. Singularity of every arrow around a
Frobenius orbit does not imply nilpotence of its product:
the matrices $\operatorname{diag}(0,1,\ldots,1)$ already show
the linear-algebra distinction. This matrix observation is a
logical comparison, not a claimed geometric counterexample.

The underlying obstruction to a decreasing-rank induction is
also direct. For any allowed degree-$d$ pullback, trace gives
\[
 H^1(X,\mathcal O_X)\ \lhook\joinrel\longrightarrow\
 H^1(Y,\mathcal O_Y),\qquad
                 d^{-1}\operatorname{Tr}\circ a^*=1.
 \tag{PI.6}
\]
These maps commute with Frobenius. Hence an existing stable
Frobenius subspace persists under further allowed covers.
A fresh nonordinary representation can add a defect in a
larger space, but these cited results supply no decreasing
stable dimension or zero stable part of a new primitive factor.

Finally, a quotient obtained only from a subgroup of the
prime-to-$p$ deck group cannot supply an elliptic intermediate
curve over the original base: every such intermediate cover
$Z\to X$ is etale and
$g(Z)=1+[Z:X](g(X)-1)\ge2$.
An additional non-deck quotient or correspondence would require
new geometric input. This records the exact limit of the
pullback-iteration route, not a nonexistence theorem for fresh
witnesses upstairs and not a universal main-route advance.
